% Pista geo-23j-q6 · Geometria
% Procedència: PAU Catalunya 2023 · Convocatòria de juny · Sèrie 1 · Qüestió 6
\documentclass[11pt,a4paper]{article}
\usepackage{pista}
\pistainfo{Catalunya 2023}{Convocatòria de juny}{Sèrie 1}

\begin{document}

\section*{\textcolor{blauPista}{Qüestió 6}}

\noindent Siguin els plans $\pi_{1}$ i $\pi_{2}$, determinats respectivament per les equacions $\pi_{1}: x + y = 3$ i $\pi_{2}: x - z = -2$.

\subsection*{\textit{a)} \normalfont Trobeu l'equació general ($Ax + By + Cz + D = 0$) del pla $\pi_{3}$, que és perpendicular a $\pi_{1}$ i $\pi_{2}$, i que passa pel punt $P = (4, 1, 2)$. \hfill\textnormal{[0,75~punts]}}

\pista{Dos plans són perpendiculars quan ho són els seus vectors normals.
\begin{itemize}
\item El vector normal $\vec{n}_{3}$ de $\pi_{3}$ ha de ser perpendicular alhora a $\vec{n}_{1}$ i a $\vec{n}_{2}$. Quina operació entre dos vectors et dona directament un vector perpendicular als dos?
\item Amb $\vec{n}_{3} = (A, B, C)$, escriu $\pi_{3}: Ax + By + Cz + D = 0$.
\item Troba $D$ substituint-hi les coordenades de $P$.
\end{itemize}}

\subsection*{\textit{b)} \normalfont Sigui $r$ la recta d'intersecció de $\pi_{1}$ i $\pi_{2}$. Calculeu l'equació vectorial de la recta $r$. \hfill\textnormal{[0,75~punts]}}

\pista{La recta $r$ és continguda alhora en tots dos plans.
\begin{itemize}
\item \textit{Direcció.} Ha de ser perpendicular alhora a $\vec{n}_{1}$ i a $\vec{n}_{2}$: és la mateixa que has trobat a l'apartat \textit{a)} per a $\vec{n}_{3}$. Ja la tens!
\item \textit{Un punt.} Troba una solució del sistema format per les equacions dels dos plans: per exemple, fes $z = 0$ i resol per a $x$ i $y$.
\end{itemize}
Amb un punt i un vector director, l'equació vectorial és immediata.}

\subsection*{\textit{c)} \normalfont Calculeu el punt $Q$ de la recta $r$ que és més a prop del punt $P$. \hfill\textnormal{[1~punt]}}

\pista{El punt de $r$ més proper a $P$ és el peu de la perpendicular des de $P$ a la recta.

\medskip
\textbf{Drecera:} el pla $\pi_{3}$ de l'apartat \textit{a)} passa per $P$ i és perpendicular a $r$. Per tant, $Q$ és la intersecció de $r$ amb $\pi_{3}$.
\begin{itemize}
\item Escriu un punt genèric $Q$ de $r$ amb el paràmetre $t$ (equació vectorial de l'apartat \textit{b)}).
\item Escriu $\vec{PQ}$ en funció de $t$ i imposa que $Q$ sigui de $\pi_{3}$; com que $\vec{v}_{r}$ és normal a $\pi_{3}$, és el mateix que
\[
\vec{PQ} \cdot \vec{v}_{r} = 0.
\]
\item Resol l'equació en $t$ i substitueix el valor a la parametrització de $r$ per obtenir les coordenades de $Q$.
\end{itemize}}

\end{document}
